\section{Introduction}\label{sec:intro}
What does a reported security-agent attack success establish? The program
must satisfy the security predicate, admission must follow the task contract,
and evidence must bind the result to the credited execution. Each requirement
changes the interpretation of success.

ACTS-Attack investigates these requirements through complementary
cryptographic experiments. Reserved secrets and public-key relations anchor
candidate correctness. The study supplies three empirical separations.

First, admission changes a contrast with recovery fixed. On 128 matched L3
instances, the base64--hexadecimal contrast is $-0.8$ percentage points under
the declared contract and $-32.8$ under a filter rejecting prompt-permitted
imports. Unused \texttt{hashlib} imports account for $-28.9$ points of the movement;
output-preserving cleanup changes the filtered contrast to $+3.9$ points.

Second, supplied-operation feasibility (12/12), agent construction (V4 8/12),
and target consequence are distinct observations. Other cohorts locate
observation/parameter failures and demonstrate supplied-chain payload release.
Third, E5 retains a correct candidate with unresolved acquisition;
constructed bundles test bindings that prevent unsupported attribution.

Prior ACTS measures cipher identification \cite{acts}; AICrypto and
CryptanalysisBench evaluate knowledge, exploitation, and cryptanalysis
\cite{aicrypto,cryptanalysisbench}. ACTS-Attack examines the
\emph{exploit-construction gap}: translating disclosed weaknesses and
observations into recovery code, while distinguishing its cryptographic
result from admission and attribution.

These experiments investigate different links in an attack claim, rather
than estimate one pooled capability or replicate one effect.
Section~\ref{sec:external-coverage} distinguishes their coverage across
interfaces, within-family transfer, and attribution implementations.

\subsection{Research Questions}
\textbf{RQ1 (Construction and admission).} Under a fixed recovery contract,
what changes with representation, and what changes when only source
admission is regraded?

\textbf{RQ2 (Recovery to consequence).} Which observations, implementation
steps, and target predicates connect a supplied weakness to an accepted
recovery or security consequence?

\textbf{RQ3 (Supported attribution).} Which authenticated bindings distinguish
a correct candidate from evidence of the claimed execution, and how does an
explicit attribution policy behave when those bindings are defective?

\subsection{Contributions}
(1)~\textbf{Controlled construction and admission measurements.}
Matched representation trials, independent cryptographic checks, and
same-program regrading separate agent output from evaluator admission.
Source analysis and cleanup explain the large local admission distortion;
equivalence remains specific to primary L3 under the tested contract.

(2)~\textbf{An empirical account of recovery requirements and target effects.}
Prospective trajectories, explicit/reduced observations, and target controls
identify successful constructions and concrete failure mechanisms across
local contracts. Supplied-chain execution establishes payload release;
unmatched cohorts retain separate denominators and descriptive contrasts.

(3)~\textbf{An evaluated contract for supported execution claims.}
Constructed bundles and local ingress tests expose the difference between
correct candidates and supported attribution. Bounded mutation testing and
configured in-toto agreement examine the policy under stated trust
assumptions. A six-release audit supplies separate reconstruction context.
